\documentclass[11pt,a4paper]{article}

\usepackage[utf8]{vietnam}
\usepackage{amsmath,amssymb}
\usepackage{graphicx}
\usepackage{float}
\setkeys{Gin}{width=\linewidth,height=0.9\textheight,keepaspectratio}
\usepackage[margin=0.75in]{geometry}
\usepackage{xcolor}
\usepackage{array,longtable,booktabs,multirow}
\usepackage{enumitem}
\usepackage{listings}
\usepackage{hyperref}
\usepackage{fancyhdr}

\pagestyle{fancy}
\fancyhf{}
\fancyfoot[C]{Page \thepage}

\providecommand{\tightlist}{%
    \setlength{\itemsep}{0pt}\setlength{\parskip}{0pt}}
\providecommand{\passthrough}[1]{#1}

\lstdefinelanguage{SQL}{
    keywords={SELECT,FROM,WHERE,JOIN,LEFT,RIGHT,INNER,OUTER,ON,GROUP,BY,ORDER,ASC,DESC,COUNT,AVG,MAX,MIN,SUM,AS,AND,OR,NOT,NULL,IS,EXISTS,WITH,RANK,OVER,PARTITION,CASE,WHEN,THEN,ELSE,END,CREATE,TABLE,PRIMARY,KEY,FOREIGN,REFERENCES,INSERT,INTO,VALUES,UPDATE,DELETE,HAVING,DISTINCT},
    sensitive=false,
    comment=[l]{--},
    morecomment=[s]{/*}{*/},
    morestring=[b]',
    morestring=[b]"
}

\lstset{
    basicstyle=\ttfamily\small,
    keywordstyle=\color{blue}\bfseries,
    commentstyle=\color{gray},
    stringstyle=\color{red!70!black},
    showstringspaces=false,
    columns=fullflexible,
    frame=single,
    breaklines=true,
    tabsize=2,
    postbreak=\mbox{\textcolor{red}{$\hookrightarrow$}\space}
}

\begin{document}

\section{Kết nối MySQL Server bằng VS Code và sửa lỗi thường
gặp}\label{kux1ebft-nux1ed1i-mysql-server-bux1eb1ng-vs-code-vuxe0-sux1eeda-lux1ed7i-thux1b0ux1eddng-gux1eb7p}

Nguồn:
https://www.geeksforgeeks.org/sql/how-to-connect-to-mysql-server-using-vs-code-and-fix-errors/

\subsection{Mục tiêu}\label{mux1ee5c-tiuxeau}

\begin{itemize}
\tightlist
\item
  Kết nối MySQL từ Visual Studio Code
\item
  Xử lý lỗi xác thực và kết nối thường gặp
\end{itemize}

\subsection{Yêu cầu}\label{yuxeau-cux1ea7u}

\begin{itemize}
\tightlist
\item
  Visual Studio Code đã cài đặt
\item
  MySQL Server (mysqld) đã được cài và đang chạy
\item
  Mở rộng MySQL cho VS Code (ví dụ: MySQL Management Tool)
\end{itemize}

\subsection{Các bước kết
nối}\label{cuxe1c-bux1b0ux1edbc-kux1ebft-nux1ed1i}

\begin{enumerate}
\def\labelenumi{\arabic{enumi}.}
\tightlist
\item
  Mở Visual Studio Code.
\item
  Vào Extensions, tìm \passthrough{\lstinline!MySQL!} và cài đặt
  extension \passthrough{\lstinline!MySQL Management Tool!}.
\end{enumerate}

\begin{figure}[H]
\centering
\includegraphics[keepaspectratio]{images/vscode-mysql-images/vscode-01.jpg}
\caption{Cài extension}
\end{figure}

\begin{enumerate}
\def\labelenumi{\arabic{enumi}.}
\setcounter{enumi}{2}
\tightlist
\item
  Mở Explorer (Ctrl+Shift+E) - mục MySQL sẽ xuất hiện. Nhấn
  \passthrough{\lstinline!Add Connection!} để tạo kết nối mới.
\end{enumerate}

\begin{figure}[H]
\centering
\includegraphics[keepaspectratio]{images/vscode-mysql-images/vscode-02.jpg}
\caption{Add Connection}
\end{figure}

\begin{enumerate}
\def\labelenumi{\arabic{enumi}.}
\setcounter{enumi}{3}
\tightlist
\item
  Điền thông tin server: \passthrough{\lstinline!Host!} =
  \passthrough{\lstinline!localhost!}, \passthrough{\lstinline!User!} =
  \passthrough{\lstinline!root!} (mặc định), nhập
  \passthrough{\lstinline!Password!} và \passthrough{\lstinline!Port!} =
  \passthrough{\lstinline!3306!}.
\item
  Nếu gặp lỗi xác thực (ví dụ khi dùng MySQL 8 với plugin xác thực mới),
  ta có thể tạo user mới sử dụng phương thức xác thực cũ.
\end{enumerate}

\begin{figure}[H]
\centering
\includegraphics[keepaspectratio]{images/vscode-mysql-images/vscode-03.jpg}
\caption{Authentication error}
\end{figure}

\begin{enumerate}
\def\labelenumi{\arabic{enumi}.}
\setcounter{enumi}{5}
\tightlist
\item
  Mở thư mục cài đặt MySQL, vào \passthrough{\lstinline!bin!}, rồi mở
  \passthrough{\lstinline!cmd!} tại đó nếu cần thao tác trực tiếp.
\end{enumerate}

\begin{figure}[H]
\centering
\includegraphics[keepaspectratio]{images/vscode-mysql-images/vscode-04.jpg}
\caption{Open bin folder}
\end{figure}

\begin{enumerate}
\def\labelenumi{\arabic{enumi}.}
\setcounter{enumi}{6}
\tightlist
\item
  Kết nối bằng client nếu cần kiểm tra từ dòng lệnh:
\end{enumerate}

\begin{lstlisting}
mysql -u root -p
\end{lstlisting}

\begin{figure}[H]
\centering
\includegraphics[keepaspectratio]{images/vscode-mysql-images/vscode-06.jpg}
\caption{Client connect}
\end{figure}

\subsection{Xử lý lỗi kết
nối}\label{xux1eed-luxfd-lux1ed7i-kux1ebft-nux1ed1i}

\begin{itemize}
\tightlist
\item
  Nếu gặp
  \passthrough{\lstinline!ERROR 2003 (HY000): Can't connect to MySQL server on 'localhost:3306' (10061)!},
  mở Services và khởi động dịch vụ MySQL (ví dụ
  \passthrough{\lstinline!MYSQL80!}).
\end{itemize}

\subsection{Tạo user dùng xác thực cũ
(mysql\_native\_password)}\label{tux1ea1o-user-duxf9ng-xuxe1c-thux1ef1c-cux169-mysql_native_password}

Ví dụ tạo user \passthrough{\lstinline!sqluser!} với plugin
\passthrough{\lstinline!mysql\_native\_password!}:

\begin{lstlisting}
CREATE USER 'sqluser'@'%' IDENTIFIED WITH mysql_native_password BY 'password';
GRANT ALL PRIVILEGES ON *.* TO 'sqluser'@'%';
FLUSH PRIVILEGES;
\end{lstlisting}

\begin{figure}[H]
\centering
\includegraphics[keepaspectratio]{images/vscode-mysql-images/vscode-07.jpg}
\caption{Flush privileges}
\end{figure}

\subsection{Thay đổi kết nối trong VS
Code}\label{thay-ux111ux1ed5i-kux1ebft-nux1ed1i-trong-vs-code}

\begin{enumerate}
\def\labelenumi{\arabic{enumi}.}
\tightlist
\item
  Xoá kết nối cũ.
\end{enumerate}

\begin{figure}[H]
\centering
\includegraphics[keepaspectratio]{images/vscode-mysql-images/vscode-08.jpg}
\caption{Delete connection}
\end{figure}

\begin{enumerate}
\def\labelenumi{\arabic{enumi}.}
\setcounter{enumi}{1}
\tightlist
\item
  Thêm kết nối mới với:
\end{enumerate}

\begin{itemize}
\tightlist
\item
  Host: \passthrough{\lstinline!localhost!}
\item
  User: \passthrough{\lstinline!sqluser!}
\item
  Password: \passthrough{\lstinline!password!}
\item
  Port: \passthrough{\lstinline!3306!}
\end{itemize}

\begin{enumerate}
\def\labelenumi{\arabic{enumi}.}
\setcounter{enumi}{2}
\tightlist
\item
  Kết nối thành công, truy cập và tạo database, bảng và thực thi truy
  vấn từ VS Code.
\end{enumerate}

\begin{figure}[H]
\centering
\includegraphics[keepaspectratio]{images/vscode-mysql-images/vscode-09.jpg}
\caption{Query results}
\end{figure}

\subsection{Ví dụ tạo database và
bảng}\label{vuxed-dux1ee5-tux1ea1o-database-vuxe0-bux1ea3ng}

\begin{lstlisting}
CREATE DATABASE myrestaurant;
CREATE TABLE IF NOT EXISTS myrestaurant.users(
  id INT(10) NOT NULL PRIMARY KEY AUTO_INCREMENT,
  name VARCHAR(200) NOT NULL,
  phone VARCHAR(200),
  address VARCHAR(200),
  password VARCHAR(200) NOT NULL
);
INSERT INTO myrestaurant.users(name,phone,address,password)
VALUES('Gaurav','123456789','Mumbai,India','pass134'),('Sakshi','987654321','Chennai,India','pass456');
\end{lstlisting}

\begin{figure}[H]
\centering
\includegraphics[keepaspectratio]{images/vscode-mysql-images/vscode-11.jpg}
\caption{Refresh and results}
\end{figure}

\subsection{Ghi chú}\label{ghi-chuxfa}

\begin{itemize}
\tightlist
\item
  VS Code hỗ trợ quản lý cơ sở dữ liệu tiện lợi khi kết hợp extension
  phù hợp.
\end{itemize}

\subsection{Tham khảo}\label{tham-khux1ea3o}

\begin{itemize}
\tightlist
\item
  GeeksforGeeks - How to Connect to Mysql Server Using VS Code and Fix
  errors?
\end{itemize}

\end{document}

